% ============================================================
% 片段：s00-front   题名页、摘要、译者说明、目录
% 译者：lead   审校状态：初稿
% ============================================================

\maketitle

\thispagestyle{empty}

\begin{abstract}
在几何朗兰兹纲领的规范场论方法中，分歧（ramification）可以用“面算符”（surface operator）来描写：
面算符支撑在二维曲面上，正如威尔逊算符与 't Hooft 算符支撑在曲线上。本文描述
$\mathcal{N}=4$ 超杨–米尔斯理论中相关的面算符及其所依赖的参数，并分析 $S$-对偶如何作用在这些参数上。
随后，在紧化到一个黎曼面之后，我们证明：关于面算符的 $S$-对偶假设，会导出几何朗兰兹纲领在
驯服分歧（tame ramification）情形下的一个自然推广。该构造涉及仿射外尔群（affine Weyl group）
在带分歧的希格斯丛模空间的上同调上的作用，以及仿射辫群（affine braid group）在该空间上
$A$-膜或 $B$-膜上的作用。
\end{abstract}

\vspace{0.4em}

\noindent{\small\bfseries 译者说明（非原文内容）}

\noindent{\small
本文件是上述论文的简体中文学术译本。原著为 Sergei Gukov 与 Edward Witten 于 2006 年 12 月
提交（v1）、2007 年 10 月 8 日修订（v2）的预印本，后收入 \emph{Current Developments in Mathematics, 2006}
（International Press）。译本的处理原则如下：

\begin{itemize}[leftmargin=1.4em,itemsep=1pt,topsep=2pt]
  \item \textbf{公式编号与原文完全一致}。原文共 252 个带编号公式，编号为 (2.1)–(2.63)、(3.1)–(3.109)、
        (4.1)–(4.20)、(5.1)–(5.18)、(6.1)–(6.24)、(A.1)–(A.18)；译本用 \verb|\setcounter{equation}{...}|
        锁定同一编号，因此可与原文 PDF 逐式对照。图 1–9、表 1 的编号亦与原文一一对应。
  \item \textbf{页码锚点}。正文每个小节标题前有注释 \texttt{\%\ ====\ 原文 §X.X：PDF 第 N 页（印刷第 N-1 页）====}，
        其中 PDF 页码 = 原文印刷页码 + 1。
  \item \textbf{本节重点}（蓝色方框）为译者所加，用于课前速览；非原文内容。
  \item \textbf{译注}（橙色方框，全文连续编号）为审校批注，分三类：易错点预警、作者省略论证的点拨、
        物理图像补充；标 [原文疑误] 者为原文疑似笔误——\emph{公式一律照原文保留}，注中说明理由与建议。
        译注可整体关闭：在 \texttt{main.tex} 中把 \verb|\showyizhutrue| 改为 \verb|\showyizhufalse|。
  \item \textbf{字体近似}。原文的 Euler Script/Roman 字体以 \verb|\mathcal|、\verb|\mathfrak| 近似，符号含义不变。
  \item \textbf{图表}取自原文 PDF 精确裁切（300 dpi），图注中译并标注原文位置；参考文献保留原文著录与编号。
\end{itemize}
}

\vspace{0.8em}

\noindent{\small\bfseries 版权与分享}

\noindent{\small
原著 © 2006 Sergei Gukov and Edward Witten，经 arXiv 以“永久、非独占分发许可”
（arXiv.org perpetual, non-exclusive license）发布；该许可只授权 arXiv 分发原文，不授予第三方翻译或再发布的权利，
故原文著作权仍完全归作者。本译本属原文的演绎作品，仅供学术研究与教学使用；可原样转发本文全文
（不得改动正文、不得删除本声明），不得用于商业用途、不得未经许可改编或再发布（尤其投稿、出版、收录进教材）。
译注与“本节重点”系译者独立添加的评论，不代表原作者观点；译文讹误由译者负责。若原作者或 arXiv 提出异议，
将立即停止分发。引用请同时标注原文与本译本：
S.~Gukov and E.~Witten, \emph{Gauge Theory, Ramification, And The Geometric Langlands Program},
arXiv:hep-th/0612073；中译本（含译注），第 1.0 版，2026。
}

\clearpage

\tableofcontents

\clearpage
